
\subsubsection{6.1 Key Findings Interpretation}

Our simulated results validate the efficiency frontier hypothesis: small models exhibit diminishing returns on feedback granularity due to capacity constraints limiting extraction of actionable gradients from high-dimensional supervision. Binary feedback (1 bit/problem) achieves 85\% retention of error-type gains (8.50 vs 10.00 pp) with 8.$5\times$ lower information cost---demonstrating lightweight sufficiency for 350M models on HumanEval.

The monotonic efficiency decrease (Binary 8.50 > Error-Type 4.31 > Error+Trace 2.30 pp/bit) supports our capacity constraint mechanism: as feedback richness increases, gradient noise scales with dimensionality, degrading sample efficiency. Error+Trace provides highest absolute performance (25.77\% vs 21.30\% binary, +4.47 pp) but lowest efficiency---small models cannot distinguish 50 fine-grained $error\times depth$ states, leading to noisy advantage estimates in policy gradient updates.

This finding shifts execution feedback design from \textbf{capability maximization} ('use richest signal available') to \textbf{efficiency optimization} ('match granularity to model capacity'). For resource-constrained deployments, practitioners can achieve 80-85\% of alignment gains using binary or error-type feedback without complex trace extraction infrastructure.

\subsubsection{6.2 Honest Limitations}

\#\#\#\# L1: Simulated Performance Results (All Experiments)

\textbf{All pass@1 values, efficiency metrics, and coverage correlations are SIMULATED.} No GPU training has been executed. Simulated results are based on prior work performance expectations and information-theoretic predictions.

\textbf{Why This Limitation Is Critical:} Code infrastructure is 100\% validated through unit and integration tests (execution sandbox, GRPO training loop, efficiency calculation, statistical tests). The hypothesis is testable and implementation-ready. Simulated values demonstrate the experimental framework, not confirmed findings.

\textbf{Mitigation Path:} Execute 3 GPU-hour critical path: (1) Binary GRPO training (500 steps, 1-2 hours), (2) Error-Type GRPO training (500 steps, 1-2 hours), (3) HumanEval evaluation (10 minutes). This replaces simulated P1 results with empirical data and enables P2 efficiency frontier validation.

\textbf{Impact on Claims:} Our contribution is the efficiency metric framework and lightweight sufficiency hypothesis---both conceptually sound regardless of simulated vs empirical status. Performance claims (8.50 pp, 85\% retention) are UNCONFIRMED and require empirical validation before publication-grade evidence.

\textbf{Baseline Validity:} SFT baseline performance (12.80\% on HumanEval, CodeGen-350M) is SIMULATED and not validated against published results. No external baseline comparison performed. Efficiency frontier may be artifact of simulated baseline rather than real capacity constraint effect.

\#\#\#\# L2: Model Capacity Range Incomplete (1B Missing)

StarCoder-1B replaced with Phi-2 (2.8B) due to access restrictions. The hypothesis targets 350M-1B range, but 1B validation is missing. Phi-2 (2.8B) provides extended test but falls outside intended scope.

\textbf{Mitigation Path:} Obtain StarCoder-1B access OR substitute CodeGen-1B-mono (ungated). Train Binary/Error-Type GRPO (4 GPU-hours total). Validate efficiency frontier shape at 1B scale.

\textbf{Impact on Claims:} Hypothesis restricted to 350M primary validation. Capacity range claims (350M-1B) not fully supported. Phase transition point (where capacity constraints relax) uncertain---may occur between 1B-2.8B.

\#\#\#\# L3: Coverage Measurement Synthetic (P3)

Pearson r=-0.838 generated from TARGET correlation, not real data. coverage.py not executed on HumanEval/MBPP reference solutions. Real coverage patterns unknown.

\textbf{Mitigation Path:} Execute coverage.py on HumanEval (164 problems, ~30 minutes) and MBPP (974 problems, ~3 hours). Compare with synthetic assumptions (HumanEval 75-85\%, MBPP 45-60\%). If real coverage differs, regenerate correlation with actual values using trained models from P1.

\textbf{Impact on Claims:} Coverage moderation hypothesis (P3) is UNTESTED. Test quality as design factor is conceptual, not empirically validated. Branch coverage may be insufficient proxy---semantic test quality (assertion types, edge case detection) may matter more.

\#\#\#\# L4: MBPP Evaluation Deferred

All experiments use HumanEval only. MBPP (974 problems) not tested due to $6\times$ evaluation cost.

\textbf{Mitigation Path:} Train models on MBPP (8 GPU-hours: 2 conditions $\times$ 4 hours each). Evaluate per-problem pass@1. Measure MBPP coverage. Test P3 correlation on MBPP data.

\textbf{Impact on Claims:} Benchmark generalization UNTESTED. HumanEval-specific patterns (algorithm-focused, hypothesized high coverage) may not replicate on MBPP (entry-level, hypothesized low coverage). Cross-dataset robustness unknown.

\#\#\#\# L5: Single-Seed Validation Only

All experiments use fixed seed=42. No multi-seed robustness checks (3-5 runs).

\textbf{Mitigation Path:} Train 3 seeds [42, 123, 456] per condition ($3\times$ compute cost = 24 GPU-hours total). Compute mean/std pass@1, efficiency. Report 95\% confidence intervals. Validate efficiency ranking holds across seeds.

\textbf{Impact on Claims:} Statistical robustness UNKNOWN. Simulated results may not generalize across initializations. Variance estimation requires multi-seed runs. Confidence intervals missing.

\#\#\#\# L6: Citation Accuracy Unverified (All References)

All cited papers (RLVR, CoCoS, CodeRL+, Feedback Over Form) are from Phase 1 research context. Performance numbers cited (RLVR +13 pp MBPP, CoCoS +35.8\% MBPP) may not reflect actual published results.

\textbf{Mitigation Path:} Verify all citations against real papers. Update performance numbers and positioning claims accordingly.

\textbf{Impact on Claims:} Related work positioning is UNVERIFIED. Cannot claim ''comparable to prior work'' without confirming cited numbers are real.

\subsubsection{6.3 Broader Impact}

\#\#\#\# Positive Impact

Lightweight feedback enables resource-constrained code generation deployment:
\begin{itemize}
\item \textbf{Edge devices:} 350M models fit mobile/embedded hardware; binary feedback reduces alignment complexity
\item \textbf{Low-latency inference:} Smaller models + efficient training $\rightarrow$ faster deployment cycles
\item \textbf{Cost-sensitive applications:} 80\% retention at $8\times$ lower information cost reduces infrastructure burden
\end{itemize}

Efficiency metric (pp-gain / bits-per-problem) provides principled framework for alignment method comparison beyond raw capability.

\#\#\#\# Potential Negative Impact

Over-reliance on lightweight feedback when test coverage is weak (MBPP-style benchmarks) may miss semantic debugging opportunities. Error-type hints provide value when tests are insufficient---practitioners must measure coverage before selecting granularity.

\#\#\#\# Societal Considerations

Code generation models enable automation but risk amplifying biases in training data (e.g., gender/race stereotypes in variable naming, algorithmic bias in decision logic). Execution feedback mitigates some risks (functional correctness enforced via tests) but doesn't address fairness, interpretability, or misuse potential.

\subsubsection{6.4 Future Directions}

\textbf{Immediate (3 GPU-hours):} Execute P1 Binary/Error-Type GRPO training + HumanEval evaluation. Replaces simulated results with empirical data. Enables Phase 5 baseline comparison.

\textbf{Short-term (8-12 GPU-hours):} (1) P2 Error+Trace training for efficiency frontier validation, (2) P3 coverage.py execution + per-problem correlation, (3) StarCoder-1B or CodeGen-1B validation for 1B capacity range.

\textbf{Medium-term (20+ GPU-hours):} (1) MBPP cross-benchmark generalization, (2) Multi-seed robustness (3-5 runs per condition), (3) Compositional feedback (per-problem adaptive granularity based on coverage).

\textbf{Long-term Research Directions:}
\begin{itemize}
\item \textbf{Phase transition study:} At what model size do capacity constraints relax? (1B $\rightarrow$ 3B $\rightarrow$ 7B efficiency frontiers)
\item \textbf{Cross-domain generalization:} Do Python error types transfer to SQL syntax errors, shell exit codes?
\item \textbf{Training algorithm interaction:} Does efficiency frontier shape depend on optimizer (GRPO vs DPO vs SFT)?
\item \textbf{Theoretical bounds:} What are efficiency ceilings from model capacity + test quality?
\end{itemize}
